\chapter{Brukaren og skjermen}

\begin{motto}
Faget fann endeleg ein objektiv grunn å stå på: brukaren.\\ Det gløymde å spørje kven brukaren stod til teneste for.
\end{motto}

Då forma flytta inn i datamaskinen, fekk faget noko det aldri hadde hatt: ein grunn som såg ut til å vere objektiv. Heile boka har vist eit fag som leitar etter ein målestokk og berre finn smak forkledd som naudsyn eller meining utan domstol. No, på terskelen til den digitale alderen, dukka det opp ein figur som lova å avgjere alt det smaken ikkje kunne avgjere: \emph{brukaren}. Var ein form god? Spør ikkje designaren, spør ikkje kritikaren, spør brukaren — sjå om han greier å bruke henne. Verkar grensesnittet? Test det. Tel feila, mål tida, observer kvar folk går seg vill. For fyrste gong såg det ut til at faget kunne grunngje vala sine i noko utanfor seg sjølv, i målbar menneskeleg åtferd. Brukarsentreringa kom som ei frigjering frå det vilkårlege, og ho bar faget gjennom tre tiår med eksplosiv vekst. Dette kapitlet handlar om kvifor den objektive grunnen ikkje var objektiv, og kvifor brukaren — nett då han skulle redde faget frå godtykket — vart den mest verknadsfulle konstruksjonen faget nokon gong laga på vegner av nokon andre enn seg sjølv.

\section{Affordans, eller forma som lovnad}

Det er verdt å sjå kvifor brukaren vart nett no, og ikkje før. Så lenge faget forma fysiske gjenstandar, var bruken stort sett synleg i forma sjølv: ein hammar er forma som det å hamre, ein kopp som det å drikke, og hundreår med handverk hadde slipt forma mot handa til ein finst kunne ein lese bruken rett av tingen. Datamaskinen braut dette. Skjermen er ein flate som kan visa kva som helst og som ikkje liknar noko av det ho gjer; ein knapp på skjermen er ikkje ein knapp, han er eit bilete av ein knapp, og ingen tusenårig handverkstradisjon hadde slipt fram korleis han skulle sjå ut. Med eitt stod formgjevaren framfor eit objekt utan iboande form, der tilhøvet mellom det forma viste og det ho gjorde, var heilt opp til designaren å finne på. Difor måtte brukaren plutseleg konsulterast: når forma ikkje lenger sa seg sjølv, var den einaste prøva om nokon greidde å bruke henne. Skjermen skapte brukaren som problem, og brukaren vart faget si redning frå eit tomrom skjermen sjølv hadde opna.

Mannen som meir enn nokon lærte faget å sjå brukaren, var ein kognisjonspsykolog som kom til design utanfrå, slik Simon hadde gjort, og som — også her finst eit ekko — leverte ei innsikt faget tok imot med begge hender og forenkla nesten med det same. Donald Norman vende blikket bort frå gjenstanden og mot mennesket framføre han, og spurde: kvifor er somme ting umoglege å bruke, medan andre er sjølvforklarande \autocite{norman2013}? Svaret hans samla seg kring eit omgrep han lånte og omforma: \emph{affordans}. Ein dørhank affordar drag; ei plate affordar dytt; ein knapp affordar trykk. Ei god form gjer sine eigne bruksmoglegheiter synlege, slik at brukaren skjønar kva han skal gjere utan instruks, og ei dårleg form skjuler dei, slik at han må gisse, prøve og kjenne seg dum. Feilen, sa Norman, ligg som regel ikkje hjå brukaren men hjå designaren; den dørtypen folk dyttar når dei skulle dra, beviser ikkje at folk er dumme, men at forma lyg om kva ho vil ha gjort.

Lat oss merke oss noko ved sjølve omgrepet, for det røper meir enn det fyrst ser ut. Affordans var ikkje Normans eige; han lånte det frå sansepsykologien, der det tydde dei handlingsmoglegheitene ein omgjevnad faktisk byd ein organisme, uavhengig av om organismen ser dei. Norman bøygde det til å tyde dei moglegheitene forma \emph{viser fram}, og måtte seinare sjølv korrigere seg og skilje mellom kva ei form let deg gjere og kva ho \emph{signaliserer} at ho let deg gjere. Den korrigeringa er meir avslørande enn ho ser ut. For ho seier at det avgjerande i digital design ikkje er kva som er sant om forma, men kva forma kommuniserer — og dermed er vi, utan at faget heilt la merke til det, attende ved den semantiske vendinga frå førre kapittel: forma som teikn, forma som ytring. Berre at no har ytringa endeleg fått ein målestokk ho mangla då ho gjaldt Memphis-hyller, og målestokken er ikkje sanning, men effekt: virkar signalet, gjer brukaren det forma ber han om? Brukbarheit er den semantiske vendinga utstyrt med eit lab-instrument. Det høyrest ut som framsteg, og på eitt vis er det òg det. Men sjå kva slags målestokk det er: han måler om forma får brukaren til å gjere noko, ikkje om det brukaren vert fått til å gjere, er til hans gagn.

Dette var meir enn ein nyttig observasjon; det var eit heilt nytt grunnlag for å felle dom. Med affordans og brukarsentrering fekk faget for fyrste gong noko som likna ein etterprøvbar standard. Ein form var god dersom folk faktisk greidde å bruke henne, dårleg dersom dei ikkje gjorde det, og dette var ikkje smak — det var observerbart. Norman gjorde brukaren til den endelege dommaren, og lova med det å lyfte faget ut av den evige usemja om kva som var fin form og inn i ein verd der spørsmål kunne avgjerast empirisk. Sett opp mot funksjonens tomme formel og meiningas manglande domstol, var dette eit verkeleg framsteg. Faget hadde endeleg noko utanfor seg sjølv å vise til. Spørsmålet, som alltid i denne boka, er kva det ikkje hadde — og kva «utanfor seg sjølv» eigentleg viste til.

\section{Brukartesting og kunsten å måle det målbare}

Frå denne grunnen voks det fram ein heil disiplin med metodar, og det er verdt å gje metodane si ros før kritikken, for dei verka. Jakob Nielsen gjorde brukartesting til ein systematisk praksis og brukbarheit til ein målbar storleik \autocite{nielsen1993}. Han viste at ein ikkje trong store, dyre studiar for å finne dei fleste feila i eit grensesnitt; ein handfull testpersonar som freista å gjere verkelege oppgåver, observerte medan dei sleit, avslørte mesteparten av det som var gale. Han sette opp prinsipp for kva som gjer eit grensesnitt brukbart, måtar å mæle kor lang tid ei oppgåve tok, kor mange feil folk gjorde, kor nøgde dei var. For fyrste gong hadde designarbeid noko som likna tal å vise til, og tal er det språket organisasjonar høyrer. Dette var ein del av kvifor faget endeleg fekk ein plass ved bordet i teknologiselskapa: det kunne dokumentere at arbeidet verka.

Men sjå nøye på kva som faktisk vart målt, for her fell masken. Brukartesting måler \emph{effektivitet}: kor raskt, kor feilfritt, kor friksjonslaust ein brukar når eit mål. Han måler ikkje om målet er verdt å nå. Metoden tek oppgåva som gjeven — bestill varen, registrer kontoen, fullfør kjøpet — og spør berre om vegen dit er glatt nok. Spørsmålet om brukaren \emph{burde} føre vegen, om det han vert leidd til er til hans eige gagn eller mot det, fell utanfor instrumentet med full vitskapleg legitimitet, fordi instrumentet ikkje er bygd for å stille det. Eit grensesnitt som gjer det lett å seie ja og vanskeleg å seie nei, scorar høgt på brukbarheit. Eit grensesnitt som gjer det lett å melde seg på og vanskeleg å melde seg av, er, mælt med Nielsens tal, godt design. Brukbarheit er ein reell og verdfull storleik, men han er ein storleik om midlar, ikkje om mål, og faget tok han som om han var begge. Det forveksla det å måle om noko fungerer godt med det å vite om noko burde lagast. Det er den same forvekslinga vi har sett heile vegen, no kledd i empiriens autoritet: tal gjorde godtykket usynleg, slik funksjonen hadde gjort smaken usynleg.

Norman sjølv merka etter kvart at den målbare brukbarheita ikkje fanga alt som tel, og utvida ramma til å femne den emosjonelle og opplevde sida ved ein gjenstand — gleda, tilliten, det vakre \autocite{norman2004emotional}. Det var ei ærleg justering, og ho viser ein redeleg tenkjar. Men legg merke til kva utvidinga ikkje gjorde. Ho la ein dimensjon til på effektivitetssida — no skulle forma ikkje berre vere lett å bruke, men kjennast god å bruke — og gjorde med det opplevinga til endå ein storleik å optimalisere mot. Det ho ikkje la til, var spørsmålet om føremålet. Å gjere ein prosess både effektiv og emosjonelt tilfredsstillande er framleis å gjere ein prosess betre langs den aksen oppdraget peika; det seier framleis ingenting om aksen var rett. Ein gledeleg, friksjonsfri veg mot eit mål som skader brukaren, er, med det utvida apparatet, eit endå betre design enn ein berre effektiv veg dit. Faget gjorde opplevinga målbar og let framleis det einaste umålbare — om noko burde lagast — ligge urørt.

\section{Brukaren som konstruksjon}

No til det avgjerande, og det vert tydelegast i ei oppfinning som var meint som eit framsteg for brukaren og som avslører heile problemet. Alan Cooper ville ha designarane til å slutte å byggje for seg sjølve og byrje å byggje for verkelege folk, og gav dei eit verktøy: \emph{personas} \autocite{cooper1999inmates}. I staden for å designe for ein vag «brukar», skulle ein konstruere ein konkret, namngjeven fiktiv person med mål, vanar og behov, og designe for han. Det var velmeint, og det var nyttig; det disiplinerte designarbeidet og heldt sjølvopptattheita i sjakk. Men legg merke til verbet: personaen vert \emph{konstruert}. Han er ikkje funnen i verda; han er laga, og han vert laga av nokon, til eit føremål, innanfor eit oppdrag. Og her sprekk den objektive grunnen. «Brukaren» som faget gjorde til dommar, var aldri ein gjeven, nøytral storleik ein kunne lese fasiten av. Han var ein figur faget sjølv teikna — og når faget arbeider på oppdrag, vert figuren teikna i oppdragsgjevarens bilete.

Sjå kva som følgjer. Personaen får dei måla som tener forretninga: brukaren «vil» fullføre kjøpet, «ynskjer» å bli verande, «set pris på» tilrådingar. Behova som ikkje tener oppdragsgjevaren — brukarens ynske om å bruke mindre tid, kjøpe mindre, sleppe unna — vert sjeldan skrivne inn i personaen, for ingen lønner ein designar for å oppfylle dei. «Brukarsentrering» vart slik, i praksis, ikkje ei sentrering kring den verkelege personen og hans heile interesse, men kring ein konstruert person hans interesser var trimma til å samanfalle med kunden sine. Faget sa det tente brukaren, og meinte det oppriktig; men brukaren det tente var ein det hadde laga til å ville det oppdraget kravde. Den objektive grunnen var ein konstruksjon, og konstruksjonen stod til teneste for den som betalte. Dette er den same dobbeltbokføringa boka har følgt heile vegen, no i si mest sofistikerte form: faget hevdar å grunngje vala sine i brukaren, men brukaren er sjølv eit av vala, gjort på vegner av nokon andre.

\section{Stemmene som såg det komme}

Faget vart ikkje åtvara for lite. Samstundes som brukarsentreringa voks, fanst det innanfor og kring feltet ei rekkje tenkjarar som peika nett på det blinde punktet, og det er lærerikt at faget las dei, siterte dei, og likevel bygde vidare som om dei ikkje hadde skrive. Lucy Suchman viste at menneskeleg handling ikkje følgjer planar slik datamaskinen føreset, men er \emph{situert}: vi handlar i og mot ein konkret situasjon, improviserande, og eit system som modellerer brukaren som ein plan-utførar vil støtt misforstå kva han gjer \autocite{suchman1987}. Brukaren, sa ho i praksis, er ikkje den ryddige figuren maskinen treng at han skal vere; og kvar gong faget likevel bygde for den ryddige figuren, bygde det for ein abstraksjon i staden for ein person. Terry Winograd og Fernando Flores gjekk endå djupare og spurde kva slags grunnlag design for datamaskiner i det heile burde ha; dei avviste at mennesket var ein informasjonshandsamar og kravde ei forståing av at vi alltid alt er kasta inn i ein verd av meining og forplikting, ikkje tilskodarar til han \autocite{winograd1986}. Paul Dourish samla denne tråden til ein teori om \emph{kroppsleggjord samhandling}: meininga ved eit grensesnitt oppstår i bruk, i kroppen, i den sosiale og fysiske situasjonen, ikkje i hovudet til ein designar som spesifiserte han på førehand \autocite{dourish2001}.

Desse stemmene hadde, samla, alt faget trong for å sjå kva brukarsentreringa eigentleg var. Dei sa at brukaren ikkje let seg fange i ein modell utan tap; at handlinga er situert, kroppsleg og forhandla; at å gjere brukaren til ein spesifiserbar storleik er å gjere han til noko anna enn han er. Men legg merke til kva slags kritikk dette var, og kvifor faget kunne ta imot han utan å endre seg. Det var ein kritikk av at modellen av brukaren var for fattig — at han ikkje fanga den verkelege personen godt nok. Det faget aldri tok inn over seg, og som ingen av desse heilt sa rett ut, var det hardare poenget: at modellen ikkje berre var for fattig, men \emph{interessert} — at han vart teikna av nokon, for nokon, og at spørsmålet ikkje berre var om han var sann, men kven han tente. Faget kunne difor absorbere Suchman og Dourish som ei oppmoding om rikare metode, finare etnografi, meir situert forståing, og halde fram med å selje den finare forståinga til den same oppdragsgjevaren. Kritikken gjorde modellen betre. Han gjorde han ikkje mindre interessert. Og ein betre modell av brukaren, i tenesta til den som vil styre brukaren, er ikkje ei løysing. Det er ei oppgradering.

\vignett

For no var alt på plass. Faget hadde lært å sjå brukaren, å måle han, å konstruere han, og å gjere det stadig betre. Det hadde eit grensesnitt som kommuniserte, ein målestokk som talde åtferd, og ein metode som kalla alt dette teneste. Det einaste som stod att, var det steget som snur heile maskineriet på hovudet: at brukaren sjølv vert produktet. Når det som vert seld ikkje lenger er ein ting til brukaren, men brukaren si tid, hans merksemd, hans åtferd, til nokon andre — då vert kvar einaste reiskap faget her har bygd, til eit instrument for utvinning. Brukarsentreringa snur til ekstraksjon utan at ein einaste metode treng endrast. Det er der vi går no.

\vidarelesing{\fullcite{norman2013}; \fullcite{nielsen1993}; \fullcite{cooper1999inmates}; \fullcite{suchman1987}; \fullcite{winograd1986}; \fullcite{dourish2001}; \fullcite{norman2004emotional}.}
